\section{Predicting PD-L1 status in NSCLC patients using deep learning radiomics based on CT images}

Wang et al. (2025) proposed a deep learning-based radiomics (DLR) model to predict PD-L1 expression in non-small cell lung cancer (NSCLC) using CT images \cite{Lu2025-uo}. The dataset consisted of 352 NSCLC patients from the First Affiliated Hospital of Shandong First Medical University, divided into a training cohort (n = 247) and a validation cohort (n = 105). CT scans were acquired using GE Healthcare and United Imaging systems, standardized to 1.25 mm slice thickness and stored in DICOM format. Image preprocessing included resampling, Hounsfield Unit (HU) normalization to the [-1000, 400] range, and min-max scaling. Tumor regions were manually segmented with ITK-SNAP, and radiomic features were extracted from the last convolutional layer of a ResNet-50 CNN, resized to 224 x 224 pixels. A pooling technique yielded 2048 high-dimensional features, which were reduced using the LASSO method. Seven classical ML algorithms were evaluated, including support vector machines (SVM), random forest (RF), logistic regression (LR), decision trees (DT), k-nearest neighbors (KNN), and XGBoost, all implemented in scikit-learn (Python). The SVM model achieved the highest performance, with an AUC = 0.85, sensitivity = 0.80, and specificity = 0.73 on the validation cohort. When combining DLR features with clinical variables (age, sex, BMI, serum albumin, and inflammation index), the hybrid model reached an AUC = 0.91, confirming the complementary role of clinical data in improving AI-based PD-L1 prediction. Visualization was supported by Grad-CAM, highlighting relevant image regions influencing the CNN predictions. Overall, this study demonstrated that integrating ResNet-based feature extraction with classical ML classifiers significantly enhances biomarker prediction from medical images.